\section{Computational Implementation}
\label{sec:computational_implementation}

The computational implementation is organized into separate
procedures for grid generation, initialization of the solution
fields, and enforcement of boundary conditions. This separation
allows the same computational setup to be used by all numerical
solution methods.

\subsection{Grid Generation}

The computational grid is uniform and cell-centered. The grid
generation procedure calculates the grid spacings and the
coordinates of the cell centers from the domain dimensions and
the number of physical cells in each direction.

\begin{algorithm}[H]
\caption{Grid Generation}
\label{alg:grid_generation}
\begin{algorithmic}[1]
\Require Domain dimensions $L_x,L_y$ and cell counts $n_x,n_y$
\Ensure Cell-center coordinates $x_i,y_j$ and grid spacings
$\Delta x,\Delta y$
\State $\Delta x \gets L_x/n_x$
\State $\Delta y \gets L_y/n_y$
\For{$i=0,\ldots,n_x-1$}
    \State $x_i \gets (i+\frac{1}{2})\Delta x$
\EndFor
\For{$j=0,\ldots,n_y-1$}
    \State $y_j \gets (j+\frac{1}{2})\Delta y$
\EndFor
\State \Return $x,y,\Delta x,\Delta y$
\end{algorithmic}
\end{algorithm}

\subsection{Initialization of the Solution Fields}

The initial temperature distribution is prescribed as
$T(x,y,0)=y$. The temperature is therefore initialized using
the wall-normal coordinate, while the same profile is assigned
at every streamwise location. The velocity components are
calculated from the prescribed fully developed velocity profile
and remain fixed throughout the simulation.

\begin{algorithm}[H]
\caption{Initialization of Temperature and Velocity Fields}
\label{alg:field_initialization}
\begin{algorithmic}[1]
\Require Grid coordinates $y_j$, cell counts $n_x,n_y$,
and mean velocity $\overline{u}$
\Ensure Initial temperature field $T$ and velocity fields $u,v$
\State Allocate the temperature array, including ghost cells
\For{each physical cell $(i,j)$}
    \State $T_{ij} \gets y_j$
    \State $u_{ij} \gets 6\overline{u}\,y_j(1-y_j)$
    \State $v_{ij} \gets 0$
\EndFor
\State Apply the temperature boundary conditions
\State \Return $T,u,v$
\end{algorithmic}
\end{algorithm}

\subsection{Boundary-Condition Implementation}

The temperature array includes one layer of ghost cells around
the physical domain. The function
\texttt{apply\_temperature\_boundary\_conditions} updates these
values using the relations established in
Section~\ref{subsec:ghost_cell_boundary_conditions}. The
ghost-cell temperatures are calculated from the adjacent
physical-cell temperatures and the prescribed boundary values.

\begin{algorithm}[H]
\caption{Application of Temperature Boundary Conditions}
\label{alg:temperature_boundary_conditions}
\begin{algorithmic}[1]
\Require Temperature array $T$, grid dimensions $n_x,n_y$,
grid spacing $\Delta y$, inlet profile $T_{\mathrm{in}}(y)$,
and wall temperatures $T_{\mathrm{bottom}},T_{\mathrm{top}}$
\Ensure Updated temperature array $T$
\For{each streamwise index $i$, including ghost-cell indices}
    \State $T_{i,0}\gets 2T_{\mathrm{bottom}}-T_{i,1}$
    \State $T_{i,n_y+1}\gets 2T_{\mathrm{top}}-T_{i,n_y}$
\EndFor
\For{$j=1,\ldots,n_y$}
    \State $y_j\gets (j-\frac{1}{2})\Delta y$
    \State $T_{0,j}\gets 2T_{\mathrm{in}}(y_j)-T_{1,j}$
    \State $T_{n_x+1,j}\gets T_{n_x,j}$
\EndFor
\State \Return $T$
\end{algorithmic}
\end{algorithm}

The same boundary-condition procedure is reused during the
numerical solution whenever updated temperature values require
consistent ghost-cell values. The ghost cells are not independent
unknowns; they provide the values needed to evaluate the spatial
operator near the boundaries.

\subsection{Convergence Monitoring and Diagnostics}
\label{subsec:diagnostics}

Two complementary quantities are used to assess the numerical
solutions: the normalized error relative to a reference temperature
field and the residual of the discrete steady energy equation.
The error measures the difference between two temperature fields,
whereas the residual measures how closely the computed solution
satisfies the discretized governing equation.

\subsubsection{Normalized $L_p$ Error}

The error between the numerical and reference solutions is defined
for each control volume as
\begin{equation}
E_i = T_i - T_i^{\mathrm{ref}},
\end{equation}
where $T_i$ and $T_i^{\mathrm{ref}}$ denote the numerical and
reference temperatures in control volume $i$, respectively.
A normalized, area-weighted $L_p$ error measure is defined as
\begin{equation}
L_p =
\left(
\frac{
\displaystyle\sum_{i=1}^{N_{cv}}
A_i |E_i|^p
}{
\displaystyle\sum_{i=1}^{N_{cv}} A_i
}
\right)^{1/p},
\label{eq:normalized_lp_error}
\end{equation}
where $N_{cv}$ is the number of physical control volumes,
$A_i$ is the area of control volume $i$, and $p$ specifies the
error measure. The denominator represents the total computational
domain area, so the expression measures an area-weighted average
of the error magnitude raised to the power $p$, followed by the
corresponding root.

For $p=2$, Equation~\eqref{eq:normalized_lp_error} becomes
\begin{equation}
L_2 =
\sqrt{
\frac{
\displaystyle\sum_{i=1}^{N_{cv}} A_i |E_i|^2
}{
\displaystyle\sum_{i=1}^{N_{cv}} A_i
}
}.
\label{eq:normalized_l2_error}
\end{equation}
This measure is the area-weighted root-mean-square (RMS) error.
It assigns greater weight to larger local errors because the
differences are squared before averaging.

In the present study, the computational grid is uniform, such
that every control volume has area $A_i=\Delta x\Delta y$.
The cell areas therefore cancel in
Equation~\eqref{eq:normalized_l2_error}, giving
\begin{equation}
L_2 =
\sqrt{
\frac{1}{N_{cv}}
\sum_{i=1}^{N_{cv}} |E_i|^2
}.
\label{eq:rms_error}
\end{equation}
Accordingly, the error is evaluated in the implementation as
the square root of the mean squared difference between the two
temperature arrays. Only physical-cell values are included;
ghost cells are excluded from the comparison.

\begin{algorithm}[H]
\caption{Normalized $L_2$ Error Evaluation}
\label{alg:l2_error}
\begin{algorithmic}[1]
\Require Reference temperature field $T^{\mathrm{ref}}$
and numerical temperature field $T$
\Ensure Normalized error $L_2$
\State Convert both fields to floating-point arrays
\If{$\operatorname{shape}(T^{\mathrm{ref}})
\neq \operatorname{shape}(T)$}
    \State Raise a shape-mismatch error
\EndIf
\State $E \gets T-T^{\mathrm{ref}}$
\State $L_2 \gets \sqrt{\operatorname{mean}(|E|^2)}$
\State \Return $L_2$
\end{algorithmic}
\end{algorithm}

This normalized error measure is useful for comparing temperature
fields obtained on different uniform grid resolutions because
its value is not directly increased merely by increasing the
number of control volumes. Its interpretation nevertheless depends
on the reference solution and on comparing the fields over the
same physical domain.

\subsubsection{Steady-State Residual}

The steady-state residual is evaluated using the discrete spatial
operator implemented in \texttt{build\_rhs}. For each physical
wall-normal index, the function evaluates the right-hand side
of the semi-discrete energy equation and stores the resulting
values in an auxiliary array. The residual is then extracted
over the physical computational domain, excluding ghost cells.

Writing the semi-discrete equation as
\begin{equation}
\frac{\mathrm{d}T_{ij}}{\mathrm{d}t}
=
\mathcal{R}_{ij}(T),
\end{equation}
the steady-state residual at cell $(i,j)$ is
$\mathcal{R}_{ij}(T)$, corresponding to the discrete right-hand
side evaluated using the current temperature field. The RMS
residual is defined as
\begin{equation}
R_{\mathrm{RMS}}
=
\sqrt{
\frac{1}{n_x n_y}
\sum_{i=1}^{n_x}
\sum_{j=1}^{n_y}
\left[\mathcal{R}_{ij}(T)\right]^2
},
\label{eq:steady_residual_rms}
\end{equation}
where $n_x$ and $n_y$ are the numbers of physical cells in the
streamwise and wall-normal directions, respectively.

At steady state, the time derivative vanishes and the discrete
right-hand side approaches zero throughout the computational
domain. Consequently, a decreasing RMS residual indicates
convergence toward a solution of the discrete steady energy
equation.

\begin{algorithm}[H]
\caption{Steady-State Residual Evaluation}
\label{alg:steady_state_residual}
\begin{algorithmic}[1]
\Require Mesh dimensions $n_x,n_y$, temperature field $T$,
velocity fields $u,v$, $\mathrm{Pe}$, $\mathrm{Br}$,
domain dimensions $L_x,L_y$, and convection coefficient
$\alpha_c$
\Ensure RMS residual $R_{\mathrm{RMS}}$
\State Allocate a zero-filled auxiliary array $\mathcal{R}$
\For{$j=1,\ldots,n_y$}
    \State Evaluate the discrete right-hand side using
    \texttt{build\_rhs} at wall-normal index $j$
    \State Store the returned values in column $j$ of
    $\mathcal{R}$
\EndFor
\State Extract the physical-cell residuals
$\mathcal{R}_{\mathrm{phys}}
\gets \mathcal{R}_{1:n_x,\,1:n_y}$
\State $R_{\mathrm{RMS}} \gets
\sqrt{\operatorname{mean}
\left(\mathcal{R}_{\mathrm{phys}}^2\right)}$
\State \Return $R_{\mathrm{RMS}}$
\end{algorithmic}
\end{algorithm}

The RMS residual assesses convergence of the discrete governing
equation, whereas the normalized $L_2$ error assesses agreement
with a reference temperature field. These quantities are
complementary: a small residual indicates that the discrete
equation is nearly satisfied, but it does not by itself establish
the accuracy of the numerical solution or grid independence.

